\documentclass[10pt,twocolumn]{article}
\usepackage[T1]{fontenc}
\usepackage[utf8]{inputenc}
\usepackage{lmodern}
\usepackage[margin=1.8cm,top=2.4cm,bottom=2cm,columnsep=0.6cm]{geometry}
\usepackage{fancyhdr}
\usepackage{titlesec}
\usepackage{booktabs}
\usepackage{amsmath,amssymb,amsfonts}
\usepackage{microtype}
\usepackage{xcolor}
\usepackage{mdframed}
\usepackage{parskip}
\usepackage{enumitem}
\usepackage{hyperref}

\definecolor{rulered}{RGB}{180,30,30}
\definecolor{boxgray}{RGB}{245,245,245}
\definecolor{keywordgray}{RGB}{100,100,100}

\pagestyle{fancy}
\fancyhf{}
\fancyhead[L]{\textsc{\small Claude Mag}}
\fancyhead[C]{\small\textit{Machine Learning Research Digest}}
\fancyhead[R]{\small 2026-09-28}
\fancyfoot[C]{\small\thepage}
\renewcommand{\headrulewidth}{0.8pt}

\titleformat{\section}{\normalsize\bfseries\color{rulered}}{}{0pt}{}%
  [\vspace{-4pt}\color{rulered}\rule{\columnwidth}{0.4pt}]
\titlespacing{\section}{0pt}{8pt}{4pt}

\hypersetup{colorlinks=true,urlcolor=rulered,linkcolor=black}

\begin{document}
\setlength{\parindent}{0pt}
\setlength{\parskip}{4pt}

{\small\color{keywordgray}\textbf{Keywords:}
  diffusion models \textbullet{}
  flow matching \textbullet{}
  latent space \textbullet{}
  representation learning}\par\vspace{4pt}

{\LARGE\bfseries\raggedright
  On the Diffusibility of High-Dimensional Latents}\par\vspace{4pt}

{\large\itshape\raggedright
  Reconstruction-Finetuned Encoders Collapse to Low-Dimensional
  Manifolds---Switch to Data Prediction to Fix Flow Matching}\par\vspace{6pt}

{\small\textbf{By} Chao Feng, Zhiyang Xu, Bowei Chen, Yuanjun Xiong,
  Xiyao Wang, Jui-Hsien Wang, Richard Zhang, Zhe Lin, Andrew Owens
  \& Yijun Li}\par\vspace{2pt}

{\small\color{keywordgray}arXiv:2609.28473 \textbullet{} ECCV 2026}
\par\vspace{2pt}\noindent{\color{rulered}\rule{\columnwidth}{0.6pt}}\vspace{6pt}

Reconstruction-finetuned visual encoders exhibit a sharp collapse in effective
dimensionality, leaving standard velocity prediction in flow matching to fit
orthogonal noise directions outside the signal manifold; switching to
$\mathbf{x}_0$-prediction (clean data parameterization) focuses optimization
on the true signal subspace and substantially improves generation quality in
high-dimensional Representation Autoencoder latent spaces.

\begin{mdframed}[backgroundcolor=boxgray,linecolor=rulered,linewidth=1pt,
    innerleftmargin=6pt,innerrightmargin=6pt,
    innertopmargin=6pt,innerbottommargin=6pt]
\textbf{\small AT A GLANCE}\par\vspace{4pt}
\begin{description}[leftmargin=0pt,labelsep=4pt,itemsep=2pt,parsep=0pt,
    font=\small\bfseries]
  \item[Problem:] Finetuning visual encoders for reconstruction collapses
    their effective dimensionality, making standard velocity prediction in
    flow matching inefficient because it must fit noise directions orthogonal
    to the signal manifold.
  \item[Why it matters:] Using pretrained visual encoder feature spaces as
    diffusion latents is desirable for rich semantic representations, but
    generation quality is poor without addressing this geometric mismatch.
  \item[Method:] Analyze effective dimensionality collapse in reconstruction-
    finetuned RAEs; show $\mathbf{x}_0$-prediction avoids fitting orthogonal
    noise directions.
  \item[Result:] Substantial improvement in FID and CLIP scores by switching
    parameterization without any other change.
  \item[Evidence:] CLIP ViT-L/14 and DINOv2 RAEs, ECCV 2026 evaluation.
\end{description}
\end{mdframed}

\section{The Big Picture}

Latent diffusion models (LDMs) achieve high-quality image generation by
operating in the compressed latent space of a variational autoencoder (VAE)
rather than pixel space. A natural extension is to use the feature spaces of
powerful pretrained visual encoders---CLIP, DINOv2, and their variants---as
the latent space, combining the generative power of diffusion with the rich
semantic representations learned by these encoders.

Representation Autoencoders (RAEs) equip pretrained encoders with a learned
decoder to enable image reconstruction from features. But many off-the-shelf
encoders are trained for classification or contrastive objectives and discard
fine-grained visual details. To enable high-quality reconstruction, RAEs
fine-tune the encoder on reconstruction objectives.

This finetuning introduces a subtle but consequential geometric problem:
while the nominal feature dimension remains large (768, 1024, or more),
the learned features concentrate near a much lower-dimensional subspace.
Standard flow matching with velocity prediction must fit noise directions
orthogonal to this subspace, making optimization inefficient and generation
quality poor. This paper diagnoses the problem and proposes a simple,
principled fix.

\section{Why Existing Approaches Fall Short}

\textbf{Standard VAE latents} (Stable Diffusion) are designed for diffusion:
their dimensionality is explicitly controlled and their geometry is well-
conditioned. RAE latents inherit neither property by design.

\textbf{Velocity prediction} ($v$-prediction, $\varepsilon$-prediction) in
flow matching assumes data fills the ambient space. When data concentrates
on a low-dimensional manifold, the velocity field has large orthogonal
components that carry no signal but must be learned, wasting capacity.

\textbf{Higher nominal dimensionality} makes this worse: more dimensions
implies more orthogonal directions to fit, even though no additional signal
is present.

\textbf{Dimensionality reduction} (PCA projection before diffusion) partially
addresses the problem but requires separate engineering, discards
representation structure, and introduces artifacts at reconstruction.

\section{Core Mechanism}

\textbf{Effective Dimensionality Analysis.} Feature dimensionality is
quantified via the participation ratio of the feature covariance:
\[
d_\text{eff} = \frac{\left(\sum_i \lambda_i\right)^2}{\sum_i \lambda_i^2}
\]
where $\lambda_i$ are the eigenvalues of the feature covariance matrix.
RAE features show a sharp collapse: $d_\text{eff} \ll d_\text{nominal}$---
typically 5--15\% of the nominal dimension after reconstruction finetuning.

\textbf{$\mathbf{x}_0$-Prediction.} In flow matching with linear interpolation:
\[
\mathbf{x}_t = (1-t)\,\mathbf{x}_0 + t\,\boldsymbol{\varepsilon},
\quad \boldsymbol{\varepsilon} \sim \mathcal{N}(\mathbf{0},\mathbf{I}_d)
\]

Velocity target: $\mathbf{v} = \mathbf{x}_0 - \boldsymbol{\varepsilon}$,
which has a large noise component $-\boldsymbol{\varepsilon}$ that spans the
full $d$-dimensional space, forcing the model to fit orthogonal directions.

Data target: $\mathbf{x}_0$, which lies on the $d_\text{eff}$-dimensional
signal manifold. Predicting $\mathbf{x}_0$ avoids the orthogonal noise
component entirely.

\section{Technical Formulation}

Let $\mathbf{x}_0 \in \mathbb{R}^d$ be an RAE feature with
$d_\text{eff} \ll d$, and $\boldsymbol{\varepsilon} \sim \mathcal{N}(\mathbf{0},
\mathbf{I}_d)$.

The noisy sample at time $t$:
\[
\mathbf{x}_t = (1-t)\,\mathbf{x}_0 + t\,\boldsymbol{\varepsilon}
\]

Velocity-prediction loss:
\[
\mathcal{L}_v = \mathbb{E}_{t,\mathbf{x}_0,\boldsymbol{\varepsilon}}\!\left[
  \|f_\theta(\mathbf{x}_t, t) - (\mathbf{x}_0 - \boldsymbol{\varepsilon})\|^2
\right]
\]

Data-prediction loss:
\[
\mathcal{L}_{x_0} = \mathbb{E}_{t,\mathbf{x}_0,\boldsymbol{\varepsilon}}\!\left[
  \|f_\theta(\mathbf{x}_t, t) - \mathbf{x}_0\|^2
\right]
\]

The velocity residual decomposes as:
\[
\mathbf{x}_0 - \boldsymbol{\varepsilon}
= \underbrace{\mathbf{x}_0}_{\text{signal, on manifold}}
  - \underbrace{\boldsymbol{\varepsilon}}_{\text{noise, spans } \mathbb{R}^d}
\]

Fitting $\mathbf{x}_0$ only is equivalent to solving a projection onto the
signal manifold, avoiding the orthogonal noise component. The optimization
landscape for $\mathcal{L}_{x_0}$ is lower-dimensional and better-conditioned
than for $\mathcal{L}_v$ when $d_\text{eff} \ll d$.

\section{Learning or Inference Procedure}

\textbf{Training:}
\begin{enumerate}[itemsep=2pt,parsep=0pt]
  \item Finetune the pretrained visual encoder as a Representation
    Autoencoder (RAE) for reconstruction quality.
  \item Extract RAE features for the training image set.
  \item Compute $d_\text{eff}$ to characterize the signal manifold.
  \item Train a flow-matching generator with $\mathcal{L}_{x_0}$ instead
    of $\mathcal{L}_v$.
\end{enumerate}

\textbf{Inference:}
\begin{enumerate}[itemsep=2pt,parsep=0pt]
  \item Sample $\mathbf{x}_T \sim \mathcal{N}(\mathbf{0}, \mathbf{I}_d)$.
  \item Integrate the learned flow from $t=T$ to $t=0$ using the
    $\mathbf{x}_0$ prediction.
  \item Decode the generated feature $\mathbf{x}_0$ with the RAE decoder.
\end{enumerate}

\section{What the Guarantee Says}

When $d_\text{eff} \ll d$, $\mathcal{L}_v$ is provably harder to optimize
than $\mathcal{L}_{x_0}$ because the orthogonal noise components constitute
a high-dimensional regression problem with no signal---they can only be
fitted by memorization. The data parameterization eliminates this by
construction. The benefit scales with $d/d_\text{eff}$: the greater the
effective dimensionality collapse, the larger the expected improvement.

\section{Experimental Findings}

\begin{itemize}[itemsep=2pt,parsep=0pt]
  \item Evaluated with CLIP ViT-L/14 and DINOv2 RAEs at feature dimensions
    768 and 1024
  \item $d_\text{eff}$ is 5--15\% of $d_\text{nominal}$ after reconstruction
    finetuning
  \item $\mathbf{x}_0$-prediction achieves substantially better FID compared
    to velocity prediction in the same RAE latent space
  \item CLIP similarity scores also improve, confirming benefit beyond
    distributional statistics
  \item Results are robust across encoder types and RAE decoder architectures
  \item FID gap between parameterizations correlates with $d/d_\text{eff}$,
    directly supporting the manifold mismatch hypothesis
\end{itemize}

\section{Ablations and Interpretation}

\textbf{No RAE finetuning (frozen encoder):} When the encoder is not
finetuned for reconstruction, $d_\text{eff}$ is higher and the gap between
parameterizations is smaller, confirming the mechanism.

\textbf{Alternative fix (PCA projection):} Projecting to the principal
subspace before velocity prediction partially recovers quality but is less
elegant: it requires separate PCA fitting, discards features, and underperforms
$\mathbf{x}_0$-prediction.

\textbf{Time $t$ dependence:} The difficulty of velocity prediction is
greatest at intermediate $t$ values where $\boldsymbol{\varepsilon}$ and
$\mathbf{x}_0$ contribute equally; $\mathbf{x}_0$-prediction remains
well-conditioned throughout.

\textbf{Nominal vs.\ effective dimensionality:} Higher nominal dimension is
strictly worse under velocity prediction but not under $\mathbf{x}_0$-prediction,
removing the disincentive to use richer representations.

\section*{Reference}

Chao Feng, Zhiyang Xu, Bowei Chen, Yuanjun Xiong, Xiyao Wang, Jui-Hsien Wang,
Richard Zhang, Zhe Lin, Andrew Owens, Yijun Li.
\textbf{On the Diffusibility of High-Dimensional Latents}.
arXiv:2609.28473, ECCV 2026.\\
\url{https://arxiv.org/abs/2609.28473}

\end{document}
